%% ma: L10-B15-0094
%% lop: 10
%% chuong: 6
%% bai: 15
%% dang: 10.15.4
%% loai: TN4
%% muc: NB
%% dap_an: "C"
%% nguon: "Ảnh giáo viên gửi 25/11/2026 (lô 2): Bài tập trắc nghiệm – Toán 10 Bài 15 Hàm số (sự biến thiên, đồ thị), Phần I Câu 5"
%% kien_thuc: "Lớp 10 Bài 15 (đồ thị, sự đồng biến, nghịch biến của hàm số)"
%% trang_thai: chinh-thuc
%% ghi_chu: "Đồ thị vẽ lại theo hàm y = x^3 - 3x^2 + 1 (khớp các điểm (0;1), (2;-3) của hình gốc)."
\begin{ex}
Cho hàm số có đồ thị như hình bên dưới.
\begin{center}\begin{tikzpicture}[scale=.75]
\draw[truc] (-1.2,0)--(3.6,0) node[below]{$x$};
\draw[truc] (0,-3.8)--(0,2) node[left]{$y$};
\draw[phu] (0,-3)--(2,-3)--(2,0);
\node[above] at (2,0) {\footnotesize$2$}; \node[left] at (0,-3) {\footnotesize$-3$};
\node[above right] at (0,1) {\footnotesize$1$}; \node[below left] at (0,0) {\footnotesize$O$};
\draw[thick,domain=-.75:3.2,smooth,samples=80] plot (\x,{(\x)^3-3*(\x)^2+1});
\end{tikzpicture}\end{center}
Khẳng định nào sau đây là đúng?
\choice{Hàm số nghịch biến trên khoảng $(0;3)$.}{Hàm số đồng biến trên khoảng $(-\infty;1)$.}{\True Hàm số nghịch biến trên khoảng $(0;2)$.}{Hàm số đồng biến trên khoảng $(-\infty;3)$.}
\loigiai{Trên $(0;2)$ đồ thị đi xuống nên hàm số nghịch biến trên $(0;2)$.}
\end{ex}
